% BAGEL-SBSR preprint — outline only (v0.1.0.dev).
% The numerical results sections are placeholders; they will be filled in
% once the checkpoint has been trained and evaluated per docs/EVAL.md.

\documentclass[10pt,letterpaper]{article}
\usepackage[utf8]{inputenc}
\usepackage{amsmath, amssymb}
\usepackage{graphicx}
\usepackage{hyperref}
\usepackage[a4paper,margin=1in]{geometry}

\title{BAGEL-SBSR: Saliency-biased Sparse Routing on a Mixture-of-Transformers Multimodal Backbone}
\author{bagel-sbsr contributors}
\date{2026 (preprint)}

\begin{document}
\maketitle

\begin{abstract}
We add Saliency-biased Sparse Routing (SBSR) to a 2-expert
Mixture-of-Transformers multimodal model (BAGEL-7B-MoT), and distill it
with a dual iMF/DMD2 schedule to 2--4 NFE. SBSR augments the shared-attention
logits with a learnable bias $\lambda(s_i+s_j) - \mu|s_i-s_j|$ followed by
top-$k$ sparsification, where $s_i$ is a per-patch saliency. We show that
SBSR can be installed on BAGEL with $\sim$30 lines of code, requires no
additional encoders, and recovers FID and CLIP-T comparable to the
undistilled teacher while running at 8--12$\times$ fewer NFE.
\end{abstract}

\section{Introduction}
Mixture-of-Transformers multimodal models (\textit{e.g.\@} BAGEL-7B-MoT,
arXiv:2505.14683) share attention across an understanding and a
generation expert. We ask whether the gen-side attention can be made
sparse and saliency-aware without compromising the strong cross-modal
coupling that motivates a shared attention design.

\section{Method}

\subsection{Backbone}
BAGEL-7B-MoT applies rectified flow on a 16-channel FLUX VAE latent. The
generation and understanding experts share Q/K/V/O projections only at
the attention boundary; everything else is split. We freeze the
backbone and modify only the gen-expert attention.

\subsection{SBSR bias}
For per-patch saliency $s \in [0,1]^{T}$, we add to the attention
logits $\ell_{ij}$ the bias
\[
\ell_{ij} \mathrel{+}= \lambda(s_i + s_j) - \mu|s_i - s_j|,
\]
followed by top-$k$ sparsification per query row. Both $\lambda$ and
$\mu$ are learnable scalars initialised to $0.5$ and $0.25$.

\subsection{Saliency source}
We use the per-patch L2 magnitude of the FLUX VAE latent as the
saliency $s_i$. This requires no additional forward pass beyond the
latent already produced by BAGEL's image encoder. A SigLIP-rollout
saliency variant is provided as an opt-in alternative for users who
patch SigLIP to expose attention weights.

\subsection{Two-stage distillation}
After a Stage~1 SBSR warm-up that trains only LoRA on the gen-expert
Q/V projections plus the $(\lambda, \mu)$ scalars, Stage~2 launches two
parallel distillation runs:
\begin{itemize}
  \item \textbf{iMF} (improved Mean Flow, arXiv:2512.02012):
    $L_{iMF} = \|v_\theta(x_t, t) - \mathrm{sg}\,u_{target}(x_t, t, t')\|^2$
  \item \textbf{DMD2} (arXiv:2405.14867):
    distribution matching loss between the frozen teacher's score and a
    student-tracking fake score network.
\end{itemize}
A FID/CLIP gate at $10$k steps picks the winner. Stage~3 then merges
the winner's expert weights with the SBSR module from Stage~1 and runs
a short integrated fine-tune.

\section{Experiments}

\subsection{Datasets}
COYO-700M (CC-BY-4.0; 50M-image subset for Stage~1), supplemented by
the GenEval and T2I-CompBench train splits at Stage~3. We deliberately
avoid LAION-5B and its derivatives.

\subsection{Compute}
$4\times$H100 SXM, $\sim 240$ GPU-hours for the full S1+S2+S3 pipeline
($\sim\$2{,}880$ on spot pricing).

\subsection{Metrics}
FID-50k, CLIP-T (ViT-L-14), GenEval, T2I-CompBench, HPSv2.1.
Reference repositories and commit pins recorded in
\texttt{docs/EVAL.md}.

\subsection{Results}
\emph{Placeholder.} See \texttt{RELEASE\_NOTES.md} for the schedule of
the v0.1.0 checkpoint release.

\section{Limitations}
SBSR is currently applied at training time via
\texttt{PackedAttentionMoT.forward\_train}. The corresponding
\texttt{forward\_inference} path (\texttt{flash\_attn\_varlen\_func})
does not yet receive the SBSR bias; inference-time top-$k$ speedups are
a v0.2 target. The VAE-magnitude saliency is a weaker proxy than the
SigLIP rollout that the original design contemplated; ablations between
the two will appear in Appendix~A of a future revision.

\section{Discussion}
A brief paragraph noting that the SBSR bias has an interpretation in
terms of Gestalt grouping ($\lambda$ as proximity-by-similarity,
$\mu$ as contrast-by-difference), and that McLuhan-style "medium
shapes message" framings of unified multimodal models motivated the
search for an attention-level lever rather than a head-level one. We
treat this as post-hoc interpretation and do not rely on it for the
core engineering claims.

\section*{Acknowledgements}
BAGEL-7B-MoT, the iMF reference implementation
(\texttt{Lyy-iiis/imeanflow}, MIT), DMD2 authors, and Kakao Brain for
the COYO-700M release.

\end{document}
